% =====================================================================
%  Preámbulo compartido — Trabajo de grado, Ingeniería Física, EAFIT
%  Formato exigido: página Carta, márgenes sup/izq/der 3 cm e inf 2 cm,
%  interlineado 1.5, Arial 11 pt (o equivalente), citación IEEE.
%
%  Compatible con dos motores:
%    - pdfLaTeX  -> Overleaf (usa Helvetica, clon métrico de Arial)
%    - XeLaTeX   -> Tectonic local (usa el Arial real del sistema)
% =====================================================================

% ---------- Motor y tipografía ----------
% Helvetica (clon métrico de Arial) vía NFSS: funciona igual en pdfLaTeX y en
% XeLaTeX/Tectonic, de modo que el PDF local y el de Overleaf son idénticos.
% inputenc solo aplica a pdfLaTeX; XeLaTeX ya es UTF-8 nativo.
\usepackage{iftex}
\ifPDFTeX
  \usepackage[utf8]{inputenc}
\fi
\usepackage[T1]{fontenc}
\usepackage{helvet}
\renewcommand{\familydefault}{\sfdefault}

% ---------- Idioma ----------
\usepackage[spanish,es-nodecimaldot]{babel}

% ---------- Geometría de página ----------
\usepackage[letterpaper,top=3cm,left=3cm,right=3cm,bottom=2cm]{geometry}

% ---------- Interlineado 1.5 ----------
\usepackage{setspace}
\onehalfspacing

% ---------- Paquetes de apoyo ----------
\usepackage{graphicx}
\usepackage{float}
\usepackage{array}
\usepackage{booktabs}
\usepackage{multirow}
\usepackage{enumitem}
\usepackage{amsmath,amssymb}
\usepackage{siunitx}
\usepackage{titlesec}
\usepackage{fancyhdr}
\usepackage[hidelinks]{hyperref}

% Rutas donde LaTeX busca las imágenes (relativas a la raíz del repo)
\graphicspath{{assets/figuras/}{assets/fotos/}{assets/logos/}}

\sisetup{output-decimal-marker={,},group-separator={.},group-minimum-digits=4,input-decimal-markers={.,}}

% ---------- Estilo de títulos ----------
\titleformat{\section}{\normalfont\bfseries\large}{\thesection.}{0.6em}{}
\titleformat{\subsection}{\normalfont\bfseries\normalsize}{\thesubsection}{0.6em}{}
\titlespacing*{\section}{0pt}{8pt}{3pt}
\titlespacing*{\subsection}{0pt}{5pt}{2pt}

% Espaciado de tablas y leyendas (no lo fija el formato; se ajusta para
% aprovechar el límite de páginas de contenido efectivo).
\setlength{\abovecaptionskip}{4pt}
\setlength{\belowcaptionskip}{0pt}
\setlength{\intextsep}{6pt plus 2pt minus 2pt}

% ---------- Encabezado y pie ----------
\pagestyle{fancy}
\fancyhf{}
\renewcommand{\headrulewidth}{0.4pt}
\fancyhead[L]{\footnotesize Escuela de Ciencias Aplicadas e Ingeniería}
\fancyhead[R]{\footnotesize Programa de Ingeniería Física}
\fancyfoot[C]{\footnotesize\thepage}

\setlist[itemize]{topsep=2pt,itemsep=1pt,parsep=0pt,leftmargin=1.2em}
\setlist[enumerate]{topsep=2pt,itemsep=1pt,parsep=0pt,leftmargin=1.6em}

% ---------- Atajos del proyecto ----------
\newcommand{\cop}[1]{\$\,\num{#1}}

% Datos institucionales (se pueden sobrescribir en cada main.tex)
\providecommand{\docTipo}{DOCUMENTO DE TRABAJO DE GRADO}
\providecommand{\docTitulo}{Título del documento}
\providecommand{\docFecha}{\today}
\providecommand{\docAutor}{Alexander Saldarriaga Vélez}
\providecommand{\docAsesor}{Hugo Alberto Murillo Hoyos}
\providecommand{\docCoordinador}{Elena Montilla Rosero}
